\section{Æggebekendtgørelsen}

\begin{frame}
  \vspace{3em}
  \begin{center}
    {\LARGE \textbf{"Æggebekendtgørelsen"}
    }
    \\
    \vspace{1em}
    Bekendtgørelse om æg fra vilde fugle og registrering af ægsamlinger
    \\
    \vspace{1em}
    BEK nr 936 af 27/06/2016, §§ 1 og 2
  \end{center}
\end{frame}


\begin{frame}{Bekendtgørelse om æg fra vilde fugle - Omfang (§ 1)}
\vspace{1em}

\small

\raggedright

\begin{itemize}

\itemsep0.6em % Justér afstand mellem bullets

\item Bekendtgørelsen omfatter æg fra vilde fugle, der er naturligt forekommende i den danske natur, samt æg fra arter omfattet af Fuglebeskyttelsesdirektivet (2009/147/EF) og underarter.

\item Gælder også tomme æg samt dele og produkter af æg.

\item Styrelsen for Vand- og Naturforvaltning afgør, om et æg er omfattet, og kan ved tvivl kræve anmeldelse/registrering.

\item Omfatter \textbf{ikke} æg fra landbrugsmæssigt opdræt af visse dyrearter.

\end{itemize}

\end{frame}



\begin{frame}{Bekendtgørelse om æg fra vilde fugle - Forbud og undtagelser (§ 2)}
\vspace{1em}

\small

\raggedright

\begin{itemize}

\itemsep0.6em % Justér afstand mellem bullets

\item Æg fra vilde fugle må \textbf{ikke} indsamles i naturen, opbevares, besiddes, modtages, købes, sælges, overdrages, udlånes, transporteres eller udstilles – medmindre anden lovgivning tillader det.

\item Undtagelser gælder for:
    \begin{itemize}
    \itemsep0.3em
    \item Indsamling af æg fra gråand i kunstige rugepladser
    \item Æg fra ødelagte reder af fasan, gråand eller agerhøne (til opdræt)
    \item Æg fra fasaner indfanget med særlig tilladelse fra styrelsen
    \end{itemize}

\item Den, der besidder æg fra fangenskabsopdrættede fugle af omfattede arter, skal kunne \textbf{godtgøre}, at æggene ikke stammer fra vilde fugle.

\end{itemize}

\end{frame}